\documentclass{scrartcl}
\usepackage[utf8]{inputenc}
\usepackage{hyperref}
\usepackage{url}
\usepackage{natbib}
\usepackage{graphicx}
\usepackage{float}
\usepackage{listings}
\usepackage{xcolor}
\usepackage{booktabs}
\usepackage{amsmath}


\usepackage{cleveref} % this must be the last package to be loaded

\title{\LARGE
  Odds-and-Evens Game and \\ Smart Home Alarm System
}
\subtitle{Concurrent and Distributed Programming\\a.y. 2025--2026}

\author{
    Buda Marco \and Merighi Daniele \and Turchi Jacopo
}

\date{\today}

\begin{document}

\maketitle

\section{Smart Home Alarm System}

This exercise is about the implementation of software for a hypothetical smart home alarm control system using the actor programming model.

\subsection{Identified Requirements}

The central control system is placed within the context of a network of sensors.

TODO.

\subsection{Design and Implementation Strategy}

Apache Pekko is chosen as the framework to enable development through the actor model, along with Scala as the supporting programming language. The proposed solution features three kinds of actors: the central \texttt{SmartHomeAlarmSystem}, a number of peripheral \texttt{SensorActor}s, and a \texttt{Home} actor which serves as a guardian for the other actor instances and allows for their interaction with the outside world.

All actors are implemented using the functional style (i.e. ``actor as a function'' philosophy), allowing for explicit management of the program's state.

\paragraph{Protocol} The \texttt{SmartHomeAlarmSystem} actor accepts four kinds of command messages:
\begin{itemize}
\item \texttt{Arm}, which is handled in the \texttt{disarmed} state and contains an entered PIN code along with an optional custom arming mode.
\item \texttt{Disarm}, which is handled in every state except \texttt{disarmed} and contains an entered PIN code.
\item \texttt{HandleDelayEnd}, which is used internally to signal the end of a timer for exit or entry delay.
\item \texttt{HandleSensorFiring}, which is handled in the \texttt{armed} state and specifies the source sensor, so that the system can check whether it should react depending on the current arming mode.
\end{itemize}

\end{document}